\section{The kinetic response on the selected bridge}
\label{sec:ex-selected-kinetic-response}

We supply the fine kinetic matrix bound used in
Section~\ref{sec:ex-stationary-determinant-prices}, and identify its
three derivative kernels on the selected finite periods. Let
$F_n$, $0\leq n\leq2$, be the response arrays of
Theorem~\ref{thm:ex-soft-response-mockups}, with width at most
$C\geq0$ at the Laurent weight $t_*>0$ (the constants $K$ and $\eta$
of that theorem). Retain all the theorem's hypotheses on the incoming
path and metric; in particular the response identity (SR2) is
available on the tori of the selected period sequence. Let $\Phi$ be the actual bridge
of Section~\ref{sec:ex-selected-stationary-linear-data}, with fine
differences through order one bounded by $D\mu^{1+j}$.
Fix $0<\sigma\leq1$ in its strip and require
$\mu\sigma\leq t_*$. This last condition is imposed by increasing
$L$ after $t_*$ and $\sigma$ have been chosen.

\subsection{Contracting the response arrays}
For external momenta $p_1,\ldots,p_n$ and coarse coordinate legs
$a_1,\ldots,a_n$, put $v_j=\Phi(p_j)e_{a_j}$. For an internal
input momentum $q$ the output momentum is
$q'=q+\sum_jp_j$. Define a matrix $K_n(p_1,\ldots,p_n;q)$
by its bilinear pairing with arbitrary internal vectors $u,w$:
\[
 u^T K_n(p_1,\ldots,p_n;q)w
 =\mathcal B_n(p_1,\ldots,p_n,-q',q)
                         [v_1,\ldots,v_n,u,w].
\]
Here $\mathcal B_n$ is the block-periodized multilinear form
obtained by inserting the array $F_n$ in (SR2) and summing the
base site over one fine residue block. Translation of a residue
across a block boundary carries the corresponding Bloch phase.
This defines the pairing on coordinate vectors, hence defines
the matrix uniquely. We only use momenta for which $p_j,q,q'$
belong to the width-$\sigma$ strip.

Every stitch uses words of length at most one, and its lengths obey
\begin{equation}\label{eq:ex-kinetic-stitch-count}
                  (n+2)+\sum_{j=1}^{n+2}|\ell_j|=4.
\end{equation}
This is the precise derivative count: two differences for the value,
one for the first background derivative, and none for the second.
Set $R=D+5e^4$. In normalized counting measure on a residue block,
the external leg with word $\ell_j$ has supremum norm at most
$\mu^{1+|\ell_j|}R$. A unit internal leg in the weighted Fourier
norm has normalized $L^2$ norm at most the same expression, by
the multiplier estimate of
Section~\ref{sec:ex-weighted-frame-penalty}. There are exactly
two such internal legs. H\"older's inequality therefore bounds
the normalized sum of their product and all external factors by
\[
                         \prod_{j=1}^{n+2}\mu^{1+|\ell_j|}R.
\]
Restoring the unnormalized residue sum multiplies by $L^4$.
Equation~\eqref{eq:ex-kinetic-stitch-count} gives the cancellation
\[
        L^4\prod_{j=1}^{n+2}\mu^{1+|\ell_j|}R=R^{n+2}.
\]
Thus neither a count of fine sites nor a conversion of a full
matrix to entrywise bounds introduces an extra power of $L$.

We also check complex translation factors. For a shift $z$ of a
residue leg, its Bloch extension differs from the permuted residue
vector by a phase with modulus at most
$\exp(\sigma(8+\mu|z|_1))$. To see the bound, write the translated
site as its residue plus $L$ times its block quotient. The difference
between the two residues has coordinate sum at most $8L$; the
remaining displacement contributes $|z|_1$. Multiplying the phases
of the $n+2$ legs costs at most
\[
                 e^{8\sigma(n+2)}e^{\mu\sigma|z|_1},
\]
where now $z$ is the whole array shift. The last factor is paid
by its Laurent width. Summing shifts, color addresses and stitch
addresses therefore gives
\begin{equation}\label{eq:ex-selected-kinetic-bound}
 \|K_n\|\leq A_n,
 \qquad A_n=N_n e^{8\sigma(n+2)}R^{n+2}C,
 \qquad 0\leq n\leq2,
\end{equation}
where $N_n$ is the product of the numbers of color and stitch
addresses at level $n$. These numbers depend only on the fixed
group and $n$. The norm is the operator norm between the weighted
fine space and its bilinear dual at the two indicated momenta.
Taking $A_K=A_0+A_1+A_2$ supplies the common bound required in
the determinant comparison.

Absolute Laurent convergence, uniformly on every strictly smaller
strip, also proves holomorphy of these matrices. Moving a momentum
by a reciprocal lattice vector changes its residue phases and its
block quotient phases by inverse factors, so the matrices are
periodic. This verifies the common analytic extension directly
from the arrays and the bridge.

\subsection{Identification with the actual Hessian derivatives}
Let $I_m$ be the incoming action on a fine torus, as a function of
logarithmic link coordinates, and let $\mathcal D=\mathsf D_m$ be the
split-field map of Section~\ref{sec:ex-soft-curvature-receipts}: on
every link $e$,
\[
 \mathcal D(x_+,x_-)_e
       =\log\bigl(\exp(x_{+,e})\exp(x_{-,e})\bigr).
\]
Near zero,
$I_m(\mathcal D(x_+,x_-))$ is exactly the action at the background
configuration $\exp(x_-)$ after the left increment $\exp(x_+)$. Its Hessian in $x_+$ is therefore
the fine kinetic Hessian in the left coordinates used for the
stationary determinant. Differentiating it $n$ times in the
background $x_-=Wx$, and testing with two arbitrary internal
vectors, gives precisely the mixed response with $n$ minus-species
legs followed by two plus-species legs. Analyticity permits this
reordering of derivatives. This is the species order in (SR2).

On a coarse Fourier wave, the actual installed matrix $W$ produces
the block wave with residue amplitude $\Phi(p)e_a$, by the
finite Fourier identity of
Section~\ref{sec:ex-selected-stationary-linear-data}. Insert these
waves and the two internal block waves into (SR2). The sum over
coarse base sites is zero unless the total coarse momentum vanishes;
otherwise it is the coarse site cardinality. In the latter case
the remaining residue and shift sum is exactly the defining
pairing of $K_n$ above. Fourier waves span the full finite-volume
space, and testing against every outgoing wave determines a matrix.
Thus $K_0,K_1,K_2$ are the actual value, first derivative and
second derivative kernels of the same kinetic Hessian on every
selected torus, including nonzero momentum transfer. They are
not only limiting coefficient pictures.

Finally the local bridge legs used by the frame and penalty estimates
obey their required common pointwise bound: at a residue the color
Euclidean norm is at most $\sqrt{d_1}\,D\mu$, because there are
$d_1$ one-form color coordinates. Enlarging the local price to
$\sqrt{d_1}D+e^4$ also covers the fixed coordinate legs. This price
is independent of $L$. Consequently the local estimates, the kinetic
triple and all finite-volume identifications in the determinant
comparison refer to one and the same selected bridge and path.
